\chapter{PROBLÈMES RENCONTRÉS ET SOLUTIONS PROPOSÉES}
\thispagestyle{empty}

\section{Problèmes méthodologiques}

Le problème le plus important rencontré pendant cette phase a été de nature
méthodologique. La branche d'évaluation antérieure utilisait une séparation
sur des arêtes sortantes orientées, alors que les modèles structurels, en
particulier Node2Vec, opéraient finalement sur une structure de graphe non
orientée. Cela a créé un décalage entre la définition du benchmark et la façon
dont le signal structurel était consommé par le modèle. L'effet a été
particulièrement important pour Node2Vec, dont la domination antérieure est
devenue discutable dès que la branche d'évaluation a été redessinée.

Un second problème méthodologique concernait l'équité entre familles de
modèles. Réutiliser d'anciens caches GCN et GraphSAGE aux côtés de la branche
Node2Vec révisée aurait produit une comparaison hétérogène. Ce problème a été
évité en créant des notebooks distincts de réentraînement pour la branche
structurelle révisée.

\section{Problèmes d'implémentation et d'ingénierie}

Plusieurs problèmes pratiques sont également apparus au cours de
l'implémentation :
\begin{itemize}
    \item un notebook d'évaluation est devenu mal formé au niveau JSON alors
    même que la logique du code restait pertinente,
    \item des caches obsolètes dans l'ancienne branche d'évaluation ne
    correspondaient plus à la logique corrigée,
    \item l'arrêt anticipé dans la première version des notebooks GNN révisés
    utilisait le mauvais ensemble d'arêtes et risquait donc d'introduire une
    fuite d'information depuis le test,
    \item l'application web n'exposait initialement qu'un sous-ensemble de
    modèles et devait être réalignée avec la famille expérimentale complète,
    \item la gestion des titres et l'interaction de recherche dans l'interface
    nécessitaient un nettoyage, car les identifiants d'articles Wikipedia
    contiennent des underscores alors que l'interface destinée à l'utilisateur
    doit afficher des titres naturels,
    \item la couche expérimentale de LLM local pour la génération de parcours
    d'apprentissage s'est révélée beaucoup moins stable que la couche de
    récupération ordonnée. Les sorties du petit modèle étaient souvent
    incohérentes, trop concentrées dans une seule catégorie, ou faiblement
    justifiées au niveau sémantique.
\end{itemize}

\section{Solutions adoptées ou prévues}

Les solutions adoptées pendant cette phase font désormais partie de la base du
projet :
\begin{itemize}
    \item un nouveau notebook de phase 0 crée une branche structurelle non
    orientée et exporte des caches Node2Vec sans fuite,
    \item de nouveaux notebooks GCN et GraphSAGE réentraînent les bases de
    référence GNN structurelles sur la même branche révisée,
    \item l'évaluation finale à dix modèles est centralisée dans un notebook
    dédié et mise en cache sous un espace de noms séparé,
    \item les caches obsolètes issus de branches invalides ont été supprimés ou
    isolés au lieu d'être réutilisés silencieusement,
    \item l'application web a été mise à jour pour exposer tous les modèles
    testés et offrir une expérience de recherche et de prévisualisation plus
    stable.
\end{itemize}

La couche de parcours d'apprentissage fondée sur un LLM a également dû être
repensée pendant cette phase. La première tentative demandait au modèle local
de produire en une seule fois un parcours JSON structuré sur l'ensemble du
top-20 recommandé. Cette approche était trop fragile pour un petit modèle
local et produisait fréquemment des sorties mal formées ou incomplètes.
L'implémentation actuelle utilise donc une stratégie plus simple de
classification par article en deux étapes : chaque article candidat est
d'abord placé dans un groupe large, puis raffiné vers l'une des sections
finales (fondations, cœur, avancé, adjacent). Cette solution ne résout pas
complètement le problème de qualité, mais elle produit un comportement de
prototype plus robuste que la génération monolithique initiale.

Même après cette refonte, les limites de l'organisateur local restent nettes.
Comme le modèle actuel est de petite taille, il produit parfois des
classifications sémantiquement faibles ou regroupe trop d'articles dans des
catégories similaires. Pour cette raison, le présent rapport traite la couche
de parcours d'apprentissage comme une extension expérimentale plutôt que comme
un composant validé du système de recommandation. Un modèle local plus grand,
ou un modèle suivant mieux les instructions avec une meilleure cohérence de
classification, améliorerait vraisemblablement cette partie du pipeline de
manière substantielle.

Deux directions de solution supplémentaires restent actives pour la phase
suivante :
\begin{itemize}
    \item améliorer l'organisateur de parcours d'apprentissage avec un modèle
    local plus fort et une stratégie de sectionnement plus fiable,
    \item étendre la branche révisée avec de futures expériences sur le pooling
    et sur la qualité de recommandation sans revenir à l'ancienne logique de
    séparation, qui était défectueuse.
\end{itemize}
